\subsection{Small integral annihilators and a fixed-field lattice bound}
\label{sec:small-integral-annihilator}

\paragraph{Finding A52 (small integral annihilators and a fixed-field lattice bound).}
Let $C/k$ be a smooth projective geometrically connected curve of genus
$g\ge2$ over a finite field of characteristic $p$. Let $D\to C$ be
geometrically connected and etale, geometrically Galois with an odd cyclic
group $G=\langle\sigma\rangle$ of order $n$ prime to $p$. The deck group
need not be constant over $k$. Suppose that a collision of distinct
point sections at $P,Q\in C(k)$ gives compatible lifts
$P',Q'\in D(k)$ and
\[
                  \delta=[Q'-P']\in J_D(k)[p^\infty].
\tag{IA.1}
\]
Assume that there is an integral operator $A\in\mathbf Z[G]$ such that
\[
 A\delta=0,\qquad A(\zeta_n)\ne0,\qquad
 n>(\|A\|_1-1)^2,
\tag{IA.2}
\]
where $\zeta_n$ is a primitive complex $n$th root and the norm is the
sum of the absolute values of the coefficients in the degree-$<n$
representative. Then the given geometric cover is the pullback of an
etale cyclic degree-$n$ cover between genus-one curves along a separable
map
\[
 C_{\bar k}\longrightarrow E,\qquad
 \deg(C_{\bar k}/E)=p^a=e>1,\qquad e\le g-1,
\tag{IA.3}
\]
totally ramified at $P$ and $Q$. In particular (IA.2) excludes the
collision whenever $2\le g\le p$, or whenever there is no map (IA.3)
totally ramified at $P$.

The criterion permits arbitrary integer coefficients and does not
assume that the operator is defined over $k$ or that its degree is
coprime to $n$. It is a consequence of an operator on the actual class
(IA.1), not of a freely chosen torsion module.

\begin{proof}
Choose, among all integral annihilators of $\delta$ with nonzero faithful
evaluation, an operator $B=\sum b_i\sigma^i$ of least positive coefficient
norm $d$. Thus $d\le\|A\|_1$. The greatest common divisor of its nonzero
coefficients is a power of $p$. Indeed, if a prime $\ell\ne p$ divided
every coefficient, then $B=\ell B_1$ and $\ell B_1\delta=0$ would imply
$B_1\delta=0$, since $B_1\delta$ is $p$-primary. This would contradict
minimality, and faithful evaluation would remain nonzero.

The two deck orbits of $P'$ and $Q'$ are free and disjoint. Hence a
principal function $f$ with
\[
                  \operatorname{div}(f)=B(Q'-P')
\tag{IA.4}
\]
has degree exactly $d$. It is separable: if $f$ were a $p$th power,
every $b_i$ would be divisible by $p$, and the divisor of its $p$th root
would show that $B/p$ is a smaller annihilator with nonzero faithful
evaluation. Over the perfect field $\bar k$, inseparability of a
nonconstant function is equivalent to being a $p$th power.

Put $K=\bar k(D)$, let $F$ be generated by all $G$-translates of $f$,
and let $v:D_{\bar k}\to Z$ be the corresponding map of smooth curves.
Write $e=[K:F]$. The extension is separable. Poonen's controlled-field
result, applied to the finitely many degree-$d$ generators, gives
\[
                e\mid d,\qquad g(Z)\le(d/e-1)^2
\tag{IA.5}
\]
\cite[Proposition~2.1, Lemma~2.2 and Corollary~2.3]{Poonen2007Gonality}.
Since $g(D)=n(g-1)+1$ and $n>(d-1)^2$, one has $e>1$.

The group $G$ preserves $F$, so $v$ is equivariant. We first prove
injectivity on the two relevant orbits without assuming that $G$ acts
freely everywhere on $Z$. Put $x=v(P')$, and suppose $\sigma^j x=x$.
The ramification index $r_P$ of $v$ is constant along the $P'$ orbit.
Writing $f=h\circ v$, coefficient comparison in (IA.4) gives
\[
 -b_i/r_P=\operatorname{ord}_{\sigma^i x}(h)
        =\operatorname{ord}_{\sigma^{i+j}x}(h)=-b_{i+j}/r_P
\tag{IA.6}
\]
for every $i$, including zero coefficients. Thus $\sigma^jB=B$.
Evaluation at $\zeta_n$, where $B(\zeta_n)\ne0$, forces $j=0$ modulo
$n$. The same argument applies to $v(Q')$. In particular the action of
$G$ on $Z$ is faithful, and both image orbits have exactly $n$ points.

These two image orbits cannot coincide. If they did, the coefficient
comparison used in A47, with the positive ramification indices $r_P,r_Q$,
would give, for some $j$,
\[
                         r_QB=-r_PB\sigma^j.
\tag{IA.7}
\]
Iterating this equality through the odd order $u$ of $\sigma^j$ gives
$(r_Q^u+r_P^u)B=0$ in the torsion-free integral group ring, which is
impossible. The signs here are integer divisor signs in every
characteristic.

Choose any nonzero coefficient of $B$. A fibre of $v$ through its
corresponding zero or pole of $f$ contains only points in the support of
(IA.4), because it is a zero or pole fibre of $h$. Injectivity on each
image orbit and disjointness of the two image orbits show that this
fibre consists of a single point, of ramification index $e$. The same
holds for every nonzero coefficient, and equivariance then makes all
$2n$ points in the two original orbits totally ramified. In particular
\[
                           e\mid\gcd_i(b_i).
\tag{IA.8}
\]
The preceding minimality argument now shows that $e$ is a nontrivial
power of $p$, and hence $\gcd(e,n)=1$.

This last coprimality makes the $G$-action on all of $Z$ free. Otherwise
a point stabilizer would contain an element of prime order $\ell\mid n$.
Its free action on $D$ would partition the fibre divisor into
$\ell$-element orbits of equal multiplicities, implying $\ell\mid e$.
Thus $Z\to E=Z/G$ is etale of degree $n$. Since $g(Z)<n$, the identity
$g(Z)-1=n(g(E)-1)$ gives $g(Z)=g(E)=1$.

At each of the $2n$ totally ramified points the different exponent is
at least $e$: the leading derivative of a local parameter of order
$e$ vanishes because $p\mid e$, while separability makes a later
derivative nonzero. The local different formula and Riemann--Hurwitz
therefore give
\[
                  2ne\le\deg\operatorname{Diff}(v)=2n(g-1)
\tag{IA.9}
\]
\cite[Lemma~53.12.4, Tag~\texttt{0C1F}]{StacksRH}.
This proves $e\le g-1$.

Finally, faithfulness of the action on $F$ gives $K=FK^G$ inside the
Galois extension $K/K^G$. Consequently $C_{\bar k}\to E$ has degree
$e$, and the connected etale fibre product $C_{\bar k}\times_E Z$ has
function field $K$. It is $D_{\bar k}$. Ramification is unchanged by
the etale base change, so the map is totally ramified at $P,Q$.
This proves (IA.3).
\end{proof}

\paragraph{The lattice bound in one fixed prime-power tower.}
Suppose now that the geometric deck groups form a compatible
$\mathbf Z_\ell$ tower, where $\ell$ is odd and $\ell\ne p$, over the
fixed finite field $k$. Let the fixed deck Frobenius multiplier be
$a\in\mathbf Z_\ell^\times$, assumed not to be a root of unity. Put
\[
 t=\operatorname{ord}_{\mathbf F_\ell^\times}(a),\qquad
 s=v_\ell(a^t-1),\qquad
 E=\mathbf Q(\zeta_{\ell^s})^{\langle a\rangle},\qquad
 h=[E:\mathbf Q]=\varphi(\ell^s)/t.
\tag{IA.10}
\]
Here the action in the definition of $E$ is reduction modulo $\ell^s$.
For every $r\ge s$, put $n=\ell^r$ and
$\mathcal O_r=\mathbf Z_p[Z]/\Phi_n(Z)$. The minimal deck field has
degree $f_r=t\ell^{r-s}$. Its deck action $\tau$ has invariants
\[
 R=\mathcal O_r^\tau=\mathcal O_E\otimes\mathbf Z_p,
                      \qquad \operatorname{rank}_{\mathbf Z_p}R=h.
\tag{IA.11}
\]
These identifications are compatible with the embeddings from level
$\ell^s$. They follow because $\langle a\bmod\ell^r\rangle$ contains
the whole kernel of reduction to level $\ell^s$, as in (NC.11).
Invariants of the integral cyclotomic ring are the integers of the
fixed field, and flat tensoring with $\mathbf Z_p$ preserves these
invariants.

Fix a $\mathbf Z$-basis $\theta_1,\ldots,\theta_h$ of $\mathcal O_E$.
Choose integral polynomial representatives
$\theta_i=C_i(\zeta_{\ell^s})$ and put
\[
                         C_E=\sum_{i=1}^h\|C_i\|_1.
\tag{IA.12}
\]
This constant is fixed once and for all and is independent of $r$.
Let $P_n$ be the actual primitive factor of $J_{D_r}$, and define
\[
 \begin{aligned}
 N_r&=P_n(k)[p^\infty],\\
 M_r&=\exp(N_r),\\
 W_r&=[R:\operatorname{Ann}_R(N_r)].
 \end{aligned}
\tag{IA.13}
\]
The invariant ring acts on $N_r$ by actual deck endomorphisms and
semilinear descent, as in A45. Use $M_r=W_r=1$ for the zero group.
If the cover has a rational lift of the prescribed point, then
\[
 n>\left(2\min\{M_r,C_E\lfloor W_r^{1/h}\rfloor\}-1\right)^2
\tag{IA.14}
\]
excludes a collision when $2\le g\le p$, and otherwise gives the
elliptic obstruction (IA.3).

\begin{proof}
Let $I=\mathcal O_E\cap\operatorname{Ann}_R(N_r)$. The index of $I$
in $\mathcal O_E$ is exactly $W_r$: the finite quotient is $p$-primary,
and the dense map $\mathcal O_E\to R$ is surjective modulo the open
annihilator ideal. Set $L=\lfloor W_r^{1/h}\rfloor$.
The $(L+1)^h>W_r$ elements
\[
                 \sum_{i=1}^h u_i\theta_i,\qquad 0\le u_i\le L,
\tag{IA.15}
\]
contain two in the same coset modulo $I$. Their nonzero difference
$c\in I$ has a polynomial representative at level $n$ of coefficient
norm at most $C_E L$: substitute $Z^{\ell^{r-s}}$ into the fixed
$C_i$, and use coefficient differences of absolute value at most $L$.
Because $E$ is a field, this nonzero $c$ has nonzero faithful
cyclotomic evaluation.

For any hypothetical collision put $b=e_n\delta\in N_r$, where
$e_n$ is the actual primitive projector on $p$-primary groups. This
is well defined since $p\nmid n$, and it is rational since the
integral endomorphism $ne_n$ is rational. The integral operator
\[
                  A=(1-\sigma^{\ell^{r-1}})c
\tag{IA.16}
\]
vanishes on every proper conductor, kills $b$, and thus kills all of
$\delta$. Its faithful evaluation is nonzero, and
$\|A\|_1\le2C_E L$. The alternative operator
$M_r(1-\sigma^{\ell^{r-1}})$ has norm $2M_r$ and the same properties.
Apply the preceding criterion to the better bound. This proves
(IA.14).
\end{proof}

\paragraph{What the bound does and does not supply.}
The same lattice argument for the annihilator of one actual class
$b\in N_r$ replaces $W_r$ by $\#(Rb)$. The uniform condition (IA.14)
uses the annihilator of the whole actual rational group so that it
does not depend on a hypothetical collision.
Writing $R=\prod_vR_v$, with unramified residue degrees $h_v$ and
component exponents $p^{e_v}$, gives
\[
                       W_r=p^{\sum_v h_ve_v},\qquad
                       M_r=p^{\max_v e_v}.
\tag{IA.17}
\]
Thus the lattice expression can be smaller than the maximum-exponent
expression when component exponents differ sufficiently. No example
outside the union of the previously proved A45--A47 tests is asserted
here, and no universal estimate for $W_r$ is established.

In the fully transitive case $E=\mathbf Q$, choose the basis $1$.
Then $h=C_E=1$ and $W_r=M_r$, so (IA.14) is exactly the previous
prime-power inequality $n>(2M_r-1)^2$. More generally a bounded
number of rational generators does not bound the $e_v$. The valid
module $R/p^uR$ has $W_r=p^{uh}$ and $W_r^{1/h}=p^u$, with no saving
over its exponent. This last module is an algebraic illustration,
not a claimed Jacobian-cover realization. The variable-quotient
single-packet supply of A49 does not bound the remaining component
exponents or target a prescribed tower. Hence the new sufficient
integral-lattice condition leaves an arithmetic witness-supply
requirement.
